\subhead{Abstract}
In a long compartment, smoke can descend far from the fire even when a ceiling vent captures smoke above it. We used Fire Dynamics Simulator (FDS) to study how side openings affect ceiling-jet temperatures and the exhaust flow needed to maintain a clear layer. Comparisons with propane-fire measurements showed that FDS captured the shift in peak temperature towards the closed side of the compartment, although the magnitude of the temperature error depended on door state and fire power. For a propylene fire, we replayed a search over 16 curtain and exhaust designs using Bayesian optimisation (BO) and a genetic algorithm (GA). Both found the design with the largest clearance margin earlier than the fixed search order. All 16 fell below the target, prompting changes to the exhaust and make-up air layout. In a separate replay, a physics-informed neural network (PINN) predicted minimum layer heights to prioritise the six door-state and soot-yield combinations for FDS checks. This ranking yielded no saving over the fixed order. Refined-grid FDS runs identified 1.12~m$^3$/s as the lowest-power feasible flow tested. Across all six scenarios, the layer stayed at least 55.8~mm above the 0.38~m target at all eight stations throughout the 120~s simulations.

\reportheading{1.\quad Introduction}
Room-fire models often describe the hot upper layer by a mean temperature, while tunnel-fire correlations describe how ceiling temperature decays with distance from the fire. Peng et al.\ (2021) examined a geometry that combines features of both, a long enclosure with openings along one side. Their experiments showed that the temperature near the plume depended on which doors were open, with no simple trend as the number of openings increased. For modelling, the longitudinal temperature profile therefore offers a way to test how well a model predicts both the magnitude and the location of peak heating.

In Task~1, we use Peng's propane-fire measurements to assess both features with all four doors open and with only R1 open. Doorway velocity and temperature profiles help interpret how cool air reaches the fire and where hot gas leaves the enclosure. These exchange paths also matter when ceiling fans draw replacement air through the compartment. Task~2 therefore carries both door states into the design of a smoke extraction system for a propylene fire. With end-wall make-up openings set for each door state, we seek the lowest-power tested flow that maintains a 0.38~m clear layer across all six combinations of door state and soot yield. Two independent replays ask which design to examine next in the original layout and which door and soot condition to check first at a given flow in the revised layout. The flow recommendation uses the minimum heights recorded in FDS.
